\documentclass[11pt]{article}
\usepackage{amsmath,amssymb,amsthm}
\usepackage[margin=1.2in]{geometry}
\usepackage{booktabs}
\usepackage{hyperref}

\newcommand{\Tp}{T_p}
\newcommand{\Ta}{T_a}
\newcommand{\kk}{k}
\newcommand{\tauu}{\tau}
\newcommand{\dd}{\mathrm{d}}
\DeclareMathOperator{\Cov}{Cov}
\DeclareMathOperator{\diag}{diag}

\title{Van Loan's Method for the Active-Diffusion Covariance\\
\large Why the transition covariance is a block of a matrix exponential,\\
and why the closed form had to go}
\author{}
\date{}

\begin{document}
\maketitle

\begin{abstract}
The forward process of the active diffusion model requires the transition covariance
$M(t)=\int_0^t e^{As}Qe^{A^{\!\top}\!s}\,\dd s$. This note derives Van Loan's identity,
which obtains $M(t)$ as an off-diagonal block of a single $4\times4$ matrix exponential,
explains why the previously used hand-derived closed form failed --- catastrophically at
$\kk\tauu=1$ and by up to $241\%$ at small $t$ elsewhere --- and records the numerical
verification. The essential observation is that $M(t)$ never was an integral that needed
evaluating: it is the solution of a linear ODE, and augmenting the state absorbs the
forcing term into a homogeneous system whose solution is an exponential.
\end{abstract}

\section{The quantity we need}

The active model evolves image pixels $x$ coupled to a persistent noise $\eta$,
\begin{align}
\dd x    &= (-\kk x + \eta)\,\dd t + \sqrt{2\Tp}\,\dd W_x, \label{eq:sdex}\\
\dd \eta &= -\frac{\eta}{\tauu}\,\dd t + \frac{\sqrt{2\Ta}}{\tauu}\,\dd W_\eta,
\label{eq:sdeeta}
\end{align}
which in vector form $z=(x,\eta)^{\!\top}$ reads $\dd z = Az\,\dd t + G\,\dd W$ with
\begin{equation}
A = \begin{pmatrix} -\kk & 1 \\ 0 & -w \end{pmatrix},
\qquad
Q \;=\; GG^{\!\top} \;=\; \diag\!\left(2\Tp,\; \frac{2\Ta}{\tauu^{2}}\right),
\qquad w \equiv \frac{1}{\tauu}.
\label{eq:AQ}
\end{equation}

Because the drift is linear, $z_t \mid z_0$ is exactly Gaussian,
\begin{equation}
z_t \mid z_0 \;\sim\; \mathcal{N}\!\big(e^{At}z_0,\; M(t)\big),
\qquad
M(t) \;=\; \int_0^t e^{As}\,Q\,e^{A^{\!\top}\!s}\,\dd s .
\label{eq:M}
\end{equation}

Every part of the pipeline depends on $M$:
\begin{enumerate}\itemsep2pt
\item \textbf{Forward noising.} A Cholesky factor of $M(t)$ converts white noise into the
      correlated pair $(d_x,d_\eta)$ added to the mean. This is the training data.
\item \textbf{The loss.} The denoising target is built from $M_{11},M_{12},M_{22}$ and
      $\det M$; see the companion loss note.
\item \textbf{The Tweedie correction} at the end of sampling, which evaluates
      $M(t_\varepsilon)$ at the smallest time on the grid.
\end{enumerate}
An error in $M$ therefore corrupts the data, the labels, and the final denoising step
simultaneously.

\section{The closed form, and its two failure modes}

The integral \eqref{eq:M} can be evaluated by hand. Writing $a=e^{-\kk t}$, $b=e^{-wt}$,
$T_x=\Tp$, $T_y=\Ta/\tauu^{2}$, the result previously used in the code was
\begin{align}
M_{11} &= \frac{T_x}{\kk}\left(1-a^{2}\right)
        + \frac{T_y}{\kk}\left[
            \frac{1}{w(\kk+w)}
          + \frac{4abk}{(\kk+w)(\kk-w)^{2}}
          - \frac{\kk b^{2}+w a^{2}}{w(\kk-w)^{2}}
          \right], \label{eq:oldM11}\\
M_{12} &= \frac{T_y}{w(\kk^{2}-w^{2})}
          \Big[\kk\left(1-b^{2}\right) - w\left(1+b^{2}-2ab\right)\Big],
          \label{eq:oldM12}\\
M_{22} &= \frac{T_y}{w}\left(1-b^{2}\right). \label{eq:oldM22}
\end{align}
These expressions are algebraically correct and numerically unusable.

\subsection{Failure 1: the removable singularity at $\kk\tauu=1$}

When $\kk\tauu=1$ we have $w=\kk$, so $(\kk-w)^{2}=0$ in \eqref{eq:oldM11} and
$\kk^{2}-w^{2}=0$ in \eqref{eq:oldM12}. Evaluating either expression raises a division by
zero. The limit exists --- the singularity is \emph{removable} --- but the formula cannot
reach it.

The underlying reason is structural. At $\kk\tauu=1$ the matrix $A$ becomes
\emph{defective}: it has the repeated eigenvalue $-\kk$ with only one eigenvector, i.e.\ a
single $2\times2$ Jordan block,
\begin{equation}
A = -\kk I + N,\qquad N=\begin{pmatrix}0&1\\0&0\end{pmatrix},\quad N^{2}=0
\qquad\Longrightarrow\qquad
e^{At} = e^{-\kk t}\begin{pmatrix}1 & t\\ 0 & 1\end{pmatrix}.
\label{eq:jordan}
\end{equation}
Any closed form obtained by diagonalising $A$ carries $1/(\lambda_1-\lambda_2)$ factors and
\emph{must} blow up there. This is why the $\kk\tauu=1$ arm of the $\tauu$ sweep
($\tauu=0.25$ at $\kk=4$) could not be run at all before the change.

\subsection{Failure 2: catastrophic cancellation at small $t$}

Away from the singularity the formula still fails, more insidiously. As $t\to0$,
$M_{11}\to 2\Tp t$, which is minute; but \eqref{eq:oldM11} assembles it from terms of order
$T_y/[\kk(\kk-w)^{2}]$. At $\tauu=0.2$, $\kk=4$ these individual terms are
\begin{equation}
\frac{T_y}{\kk}\cdot\frac{4abk}{(\kk+w)(\kk-w)^{2}} \approx 71,
\qquad
\frac{T_y}{\kk}\cdot\frac{\kk b^{2}+wa^{2}}{w(\kk-w)^{2}} \approx 72,
\end{equation}
whose difference must yield a contribution of order $10^{-8}$. That is nine to ten orders
of cancellation; \texttt{float32} carries about seven decimal digits, so nothing survives.

Table~\ref{tab:err} gives the measured relative error of \eqref{eq:oldM11} in
\texttt{float32} against the exact value, for the parameters actually used
($\kk=4$, $\Ta=6.4$, $\Tp=10^{-3}$).

\begin{table}[h]
\centering\small\setlength{\tabcolsep}{4.2pt}
\begin{tabular}{rr rrrrrr}
\toprule
$\tauu$ & $\kk-1/\tauu$ & $t=10^{-3}$ & $10^{-2}$ & $0.05$ & $0.2$ & $1.0$ & $2.0$\\
\midrule
$0.15$ & $-2.67$ & $2.5\times10^{-1}$ & $3.1\times10^{-3}$ & $3.3\times10^{-5}$ & $2.2\times10^{-6}$ & $7.5\times10^{-7}$ & $7.4\times10^{-7}$\\
$0.20$ & $-1.00$ & $\mathbf{2.4\times10^{0}}$ & $6.9\times10^{-3}$ & $9.0\times10^{-4}$ & $1.7\times10^{-6}$ & $4.5\times10^{-7}$ & $4.9\times10^{-7}$\\
$0.25$ & $0$ & \multicolumn{6}{c}{\emph{division by zero}}\\
$0.40$ & $+1.50$ & $5.5\times10^{-1}$ & $6.0\times10^{-3}$ & $9.8\times10^{-5}$ & $2.9\times10^{-6}$ & $1.1\times10^{-7}$ & $9.9\times10^{-8}$\\
$0.50$ & $+2.00$ & $2.3\times10^{-2}$ & $4.2\times10^{-3}$ & $6.5\times10^{-5}$ & $1.0\times10^{-7}$ & $1.8\times10^{-7}$ & $1.5\times10^{-7}$\\
$1.00$ & $+3.00$ & $5.9\times10^{-3}$ & $2.9\times10^{-4}$ & $4.1\times10^{-5}$ & $3.1\times10^{-6}$ & $1.7\times10^{-7}$ & $1.3\times10^{-7}$\\
\bottomrule
\end{tabular}
\caption{Relative error of the closed form \eqref{eq:oldM11} in \texttt{float32}.
The error grows as $|\kk-1/\tauu|\to0$. At $\tauu=0.2$ it exceeds $100\%$: the returned
$M_{11}$ bears no relation to the true value, and the resulting matrix can fail to be
positive definite, which is precisely how the Cholesky factorisation crashed. Damage is
confined to $t\lesssim0.05$ --- which is exactly where fine image detail is set.}
\label{tab:err}
\end{table}

The $+10^{-8}$ constants added to $M_{11}$ and $M_{22}$ in the original implementation were
an attempt to suppress the symptom.

\section{The trick, first in one dimension}

Consider the scalar inhomogeneous linear ODE
\begin{equation}
\dot y = ay + q,\qquad y(0)=0
\qquad\Longrightarrow\qquad
y(t) = \frac{q\left(e^{at}-1\right)}{a}.
\end{equation}
Now exponentiate the augmented matrix obtained by appending the constant $q$ as an extra
column with a zero row:
\begin{equation}
\exp\!\left[t\begin{pmatrix}a & q\\ 0 & 0\end{pmatrix}\right]
=\begin{pmatrix} e^{at} & \dfrac{q\left(e^{at}-1\right)}{a}\\[2ex] 0 & 1\end{pmatrix}.
\label{eq:scalar}
\end{equation}
The solution sits in the top-right entry. Nothing subtle occurred: appending a constant row
converts an \emph{inhomogeneous} linear ODE into a \emph{homogeneous} one, and homogeneous
linear ODEs with constant coefficients are exactly what matrix exponentials solve. The
forcing term has been absorbed into the generator.

Van Loan's identity is this same manoeuvre for matrices.

\section{Derivation}

\subsection{$M$ satisfies a Lyapunov ODE}

Differentiating \eqref{eq:M} with Leibniz' rule,
\begin{equation}
\dot M = A M + M A^{\!\top} + Q, \qquad M(0)=0.
\label{eq:lyap}
\end{equation}
This has the same shape as the scalar case: linear in the unknown, plus a constant forcing
$Q$. So we apply the same augmentation.

\subsection{The block generator}

Define the $4\times4$ matrix and its exponential,
\begin{equation}
C \;=\; \begin{pmatrix} -A & Q\\ 0 & A^{\!\top}\end{pmatrix},
\qquad
E(t) \;=\; e^{Ct}.
\label{eq:C}
\end{equation}
Since $C$ is block upper triangular, so is $E(t)$, and the diagonal blocks exponentiate
independently:
\begin{equation}
E(t) = \begin{pmatrix} E_{11}(t) & E_{12}(t)\\ 0 & E_{22}(t)\end{pmatrix}
     = \begin{pmatrix} e^{-At} & E_{12}(t)\\ 0 & e^{A^{\!\top}\!t}\end{pmatrix}.
\label{eq:blocks}
\end{equation}

\subsection{The off-diagonal block}

Read the $(1,2)$ block off $\dot E = CE$:
\begin{equation}
\dot E_{12} = -A\,E_{12} + Q\,e^{A^{\!\top}\!t},\qquad E_{12}(0)=0.
\end{equation}
Multiply by the integrating factor $e^{At}$:
\begin{equation}
e^{At}\dot E_{12} + e^{At}A\,E_{12}
\;=\;
\frac{\dd}{\dd t}\Big(e^{At}E_{12}\Big)
\;=\;
e^{At}\,Q\,e^{A^{\!\top}\!t}.
\label{eq:exact}
\end{equation}
The left-hand side is now an exact derivative, and the right-hand side is precisely the
integrand of \eqref{eq:M}. Integrating from $0$ to $t$ and using $E_{12}(0)=0$,
\begin{equation}
e^{At}E_{12}(t) \;=\; \int_0^t e^{As}Qe^{A^{\!\top}\!s}\,\dd s \;=\; M(t).
\end{equation}
Finally $E_{22}(t)=e^{A^{\!\top}\!t}$, so $E_{22}(t)^{\!\top}=e^{At}$, giving

\begin{equation}
\boxed{\;M(t) \;=\; E_{22}(t)^{\!\top}\,E_{12}(t)\;}
\label{eq:vanloan}
\end{equation}

\noindent
where $E_{12}$ and $E_{22}$ are the top-right and bottom-right $2\times2$ blocks of
$e^{Ct}$. One matrix exponential; no denominators; no case analysis.

The interpretive point: the integral was never something that had to be \emph{evaluated}.
It is a block of an exponential, because its integrand solves a linear ODE. The closed form
\eqref{eq:oldM11}--\eqref{eq:oldM22} was one particular way of unpacking that block by
hand, and a numerically poor one.

\section{Implementation}

\begin{verbatim}
def compute_covariance(self, t):
    k  = float(self.k);  w  = 1.0 / float(self.tau)
    Tx = float(self.Tp); Ty = float(self.Ta) / float(self.tau) ** 2
    A = torch.tensor([[-k, 1.0], [0.0, -w]],   dtype=torch.float64)
    Q = torch.tensor([[2*Tx, 0.0], [0.0, 2*Ty]],
                     dtype=torch.float64)
    Z = torch.zeros(2, 2, dtype=torch.float64)
    C = torch.cat([torch.cat([-A, Q], 1), torch.cat([Z, A.T], 1)], 0)
    E   = torch.linalg.matrix_exp(C * t.double().reshape(-1, 1, 1))
    cov = E[:, 2:, 2:].transpose(-1, -2) @ E[:, :2, 2:]   # E22^T E12
    cov = 0.5 * (cov + cov.transpose(-1, -2))        # symmetrise
    return cov.to(t.dtype)
\end{verbatim}

The assembly is done in \texttt{float64} and cast down once at the end. The explicit
symmetrisation removes the $O(10^{-16})$ asymmetry left by finite precision, which matters
only because Cholesky rejects non-symmetric input.

\subsection{Why this is numerically stable}

\begin{itemize}\itemsep3pt
\item \texttt{matrix\_exp} uses scaling-and-squaring with Pad\'e approximants: it evaluates
      $e^{Ct/2^{n}}$ from a series whose terms are of comparable magnitude, then squares
      $n$ times. The differences of nearly-equal exponentials that destroy
      \eqref{eq:oldM11} never form.
\item It never diagonalises $A$, so the defective case \eqref{eq:jordan} needs no special
      handling --- $e^{At}=e^{-\kk t}\big(I+Nt\big)$ emerges automatically. The
      $\kk\tauu=1$ point is not special to the algorithm.
\item Nothing in \eqref{eq:C} involves $\kk-w$, $\kk+w$, or any other denominator that can
      vanish.
\end{itemize}

\subsection{Cost}

One $4\times4$ \texttt{matrix\_exp} per batch element per step, in \texttt{float64}. It is
evaluated per time value, not per pixel, so at batch $128$ it is negligible beside the UNet
forward pass; measured epoch time was unchanged by the switch.

\section{Numerical verification}

Two checks, both against the shipped implementation.

\paragraph{The identity itself.} Comparing $E_{22}^{\!\top}E_{12}$ against direct
quadrature of \eqref{eq:M} on a $4\times10^{5}$-point grid:

\begin{center}
\begin{tabular}{llll}
\toprule
$\tauu$ & $t$ & $\max\big|\text{Van Loan}-\text{quadrature}\big|$ & $M_{11}$\\
\midrule
$0.25$ & $0.05$ & $7.1\times10^{-13}$ & $6.423485\times10^{-3}$\\
$0.25$ & $0.5$  & $2.1\times10^{-10}$ & $6.097628\times10^{-1}$\\
$0.25$ & $2.0$  & $3.4\times10^{-9}$  & $8.002369\times10^{-1}$\\
$0.50$ & $0.05$ & $4.5\times10^{-14}$ & $1.789678\times10^{-3}$\\
$0.50$ & $2.0$  & $4.3\times10^{-10}$ & $5.325359\times10^{-1}$\\
\bottomrule
\end{tabular}
\end{center}

\noindent
The residual is the trapezoid rule's discretisation error, not Van Loan's. Note that
$\tauu=0.25$ is included: the closed form does not exist there, but the block exponential
is unremarkable.

\paragraph{Structure.} At $\tauu=0.25$, $t=0.7$: $E_{11}=e^{-At}$, $E_{22}=e^{A^{\!\top}\!t}$
and the lower-left block vanishes, all to \texttt{allclose} tolerance --- confirming
\eqref{eq:blocks}.

\section{Companion fixes}

The same defect appears in the \emph{mean}. The $(1,2)$ entry of $e^{At}$ is
\begin{equation}
b(t) \;=\; \int_0^t e^{-\kk(t-s)}e^{-ws}\,\dd s
      \;=\; e^{-\kk t}\,\frac{e^{dt}-1}{d},\qquad d \equiv \kk - w,
\end{equation}
which is $0/0$ at $\kk\tauu=1$ with limit $t\,e^{-\kk t}$, consistent with
\eqref{eq:jordan}. Two changes were required:
\begin{itemize}\itemsep2pt
\item an explicit branch, \texttt{b = t*a if abs(d) < 1e-8 else a*torch.expm1(d*t)/d};
\item \texttt{expm1} in place of \texttt{exp(x)-1}. For small $x$, $e^{x}-1$ loses almost
      every significant digit to cancellation (at $x=10^{-8}$ in \texttt{float32},
      $e^{x}$ rounds to $1$ and the difference is $0$), whereas \texttt{expm1} evaluates
      the quantity directly.
\end{itemize}
The same treatment was applied at every site where these expressions appear: the training
loss, the Tweedie correction, the equilibrium transition, and both reverse-transition
routines of the SSCS sampler.

\section{Consequences for the recorded results}

Because $M$ enters the noising, the labels and the Tweedie step, all runs trained before
this change used a corrupted forward process at small $t$. Results measured on those
checkpoints remain valid statements about \emph{those models}, but they are not
reproducible from the current code, and comparisons across the boundary are confounded with
the numerical error --- which varies by arm, being largest exactly where $|\kk-1/\tauu|$ is
smallest. This is recorded per-arm in the results note.

\section{Generalisation}

Equation \eqref{eq:vanloan} makes no use of the specific $A$ and $Q$ in \eqref{eq:AQ}. It
holds for any linear SDE in any dimension. Adding a second persistent mode, giving $\kk$
matrix structure, or coupling colour channels requires only writing down the new $A$ and
$Q$; \texttt{compute\_covariance} is unchanged. That robustness to model changes, rather
than the numerics alone, is the strongest argument for the method.

\begin{thebibliography}{9}
\bibitem{vanloan1978}
C.~F.~Van Loan,
\emph{Computing integrals involving the matrix exponential},
IEEE Transactions on Automatic Control \textbf{23}(3), 395--404, 1978.
\bibitem{higham2005}
N.~J.~Higham,
\emph{The scaling and squaring method for the matrix exponential revisited},
SIAM Journal on Matrix Analysis and Applications \textbf{26}(4), 1179--1193, 2005.
\end{thebibliography}

\end{document}
